% block: multigrid-heading page:31 type:section
\subsection{{Multigrid Methods}}\label{sec:multigrid-methods}
% block: multigrid-principle page:31 type:ordinary-paragraph
{Multigrid methods fall into a class of iterative methods that have been developed for solving linear systems arising from elliptic partial differential equations. Multigrid methods find a good approximate solution of the linear systems from the information of some rough approximate solutions of smaller linear systems (arising from hierarchical coarser meshes) and an accurate approximate solution of the smallest linear systems (arising from the coarsest mesh). For obtaining the rough approximate solutions, stationary iterative methods (especially the Gauss--Seidel method) are often used, which are called smoothers. The name smoother comes from the fact that the Gauss--Seidel method damps quickly the high frequency components of the error or residuals for the linear systems, and the multigrid methods make the most of this property.}
% block: multigrid-convergence page:31 type:ordinary-paragraph
{The notable feature of the multigrid methods is that the convergence rate does not depend on the mesh size, though stationary iterative methods and Krylov subspace methods do. The multigrid methods are later extended to solving other linear systems arising from not only elliptic ones but other models.}
% block: multigrid-books page:31 type:ordinary-paragraph
{Multigrid methods are the subject of active research, and there are many excellent books on multigrid methods, some of them by:}
\begin{itemize}
\item Hackbusch \BookCite{91},
\item Wesseling \BookCite{204},
\item Bramble \BookCite{27},
\item Briggs et al. \BookCite{28},
\item Trottenberg et al. \BookCite{193},
\item Lottes \BookCite{125}.
\end{itemize}
{Wolfgang Hackbusch's book presents not only the basic concepts of multigrid methods using simple one-dimensional differential equations but also describes algorithms, their program codes, and detailed convergence analysis. James Lottes's book focuses mainly on algebraic multigrid methods.}
